\section{Joint guards and two Pell-root translations}
\label{sec:integer-new}

This section gives three explicit fixed-coefficient equations over unrestricted
integers.  The principal constructions have 26 variables and degree eight,
and 24 variables and degree ten.  A further restriction to canonical dyadic
witnesses gives 23 variables and degree sixteen.  They improve the inspected
repository's 27-variable octic and 26-variable degree-sixteen construction
\cite{source68,source71}.  No global minimum or priority claim outside the
inspected repository is made.

The starting point is the literal 38-variable quartic of source 71.  Its
fixed tuple and determinant guard are kept intact.  Two operations are
combined here: the simultaneous-positivity lemma of source 68 is applied
to the two remaining guards, and two integral translations of Pell-root
coordinates permit further substitutions without the usual degree increase.
The first operation is a synthesis of earlier results.  The paired
translations and the resulting 24-variable system are the additional
algebraic reduction.

\subsection{One guard for two strict inequalities}

Put
\[
 \operatorname{Pos}(p_1,p_2,p_3)
   =1+p_1^2+p_2^2+p_3(p_3+1).
\]
This takes only positive integer values and represents every positive
integer.  Indeed, apply the three-square theorem to $4n-3$ and separate its
one odd coordinate from its two even coordinates.  This gives
$4n-3=(2p_3+1)^2+(2p_1)^2+(2p_2)^2$.

\begin{lemma}[Nearest-square simultaneous positivity; source 68]
\label{lem:int-pair}
For any $x,y\in\Z$,
\[
 x>0\ \land\ y>0
 \quad\Longleftrightarrow\quad
 \exists t,p_1,p_2,p_3\in\Z:
 xy-(t^2-x-y+1)^2=\operatorname{Pos}(p_1,p_2,p_3).
\]
\end{lemma}
\begin{proof}
A solution makes $xy>0$, so $x,y$ have the same nonzero sign.  If
$x=-a$, $y=-b$, where $a,b\ge1$, then
$(t^2+a+b+1)^2>ab=xy$, a contradiction.

Conversely let $x,y\ge1$, put $N=x+y-1$, and take
$m=\lfloor\sqrt N\rfloor$.  Choose $t=m$ if $N-m^2\le m$, and
$t=m+1$ otherwise.  If $N$ is a square the error is zero.  Otherwise
$|t^2-N|\le m$ and $m^2<N$.  Hence
\[
 (t^2-N)^2<N\le xy,
 \qquad xy-N=(x-1)(y-1)\ge0.
\]
The remaining positive integer has a $\operatorname{Pos}$ representation.
The floor and square root are used only to choose a witness in this proof;
they are not operations in the polynomial.
\end{proof}

Two separate strict inequalities previously cost six private integer
coordinates.  Lemma~\ref{lem:int-pair} uses four.  When $x$ and $y$ are
quadratic polynomial expressions, its residual has degree at most four.
It is essential to keep both inputs quadratic while forming the octic:
eliminating $\zeta$ too early changes $\eta\zeta$ into a cubic expression.

\subsection{Fixed coefficients}

The six fixed integers are
\[
 \beta=2^{1064921}-1,\quad K_0,\quad M_C,\quad M_F,
 \qquad \sigma_F=2^{1113},\quad \sigma_C=2^{666104}.
\]
Here $(\beta,K_0,M_C,M_F)$ is the complete
\texttt{mixed-parity-unrounded-1064921} tuple of source 34, used without
mixing ports from other tuples.  In particular $M_F$ is already shifted
by $\beta$.  The archive supplies every actual integer as compressed
minimal unsigned little-endian bytes, its SHA-256 digest, and the finite
recipe rebuilding it from the literal K215 transition table and ordered
943-tile list.  Thus these are fixed numerals, not parameters to be chosen
after searching for a solution.  The coefficient proof and calibration
remain those of sources 34, 64 and 71 \cite{source64,source71}.

\begin{center}
\begin{tabular}{lrr}
\hline
Fixed numeral & Bit length & Minimal raw bytes\\
\hline
$\beta$ & 1,064,921 & 133,116\\
$K_0$ & 1,065,043 & 133,131\\
$M_C$ & 1,064,921 & 133,116\\
$M_F$ & 1,064,922 & 133,116\\
$\sigma_F$ & 1,114 & 140\\
$\sigma_C$ & 666,105 & 83,264\\
\hline
Total & & 615,883\\
\hline
\end{tabular}
\end{center}

\subsection{The octic in 26 integers}

Its independent integer coordinates are
\[
\begin{gathered}
 J,F,\alpha,z,f,h,r,v,w,\tau,\eta,\zeta,y,\gamma,o,j,\\
 E,c,a,V,I,u,t,p_1,p_2,p_3.
\end{gathered}
\]
The code names $v,w,t,p_i$ as \texttt{v\_scale}, \texttt{w\_scale},
\texttt{pair\_root}, and \texttt{guard1},\ldots,\texttt{guard3}.
Every expression in the following display is an abbreviation for a
polynomial in these 26 coordinates:
\begin{align*}
 q&=\beta J+1, & k&=\eta+\zeta, & Y&=a-E,\\
 C&=q-F-\alpha, & X&=2u^2, & H&=4a+3,\\
 \delta&=a^2+H, & T&=E(c-\eta), & S&=X+\gamma H,\\
 R&=\bigl((q-1)(q-F)+\alpha\bigr)(q^2-1)
       +(M_C+qM_F)J.&&&
\end{align*}
The two guard expressions and their combined residual are
\begin{align*}
 G&=(\sigma_F F)(\sigma_C C)
       -(q-\sigma_F F-\sigma_C C)^2,\\
 Z&=\eta\zeta,\\
 B_*&=GZ-(t^2-G-Z+1)^2-\operatorname{Pos}(p_1,p_2,p_3).
\end{align*}
Define the fourteen residuals
\begin{align}
 A_1&=X-4qw, & A_2&=a-E-q^3(2v+1),\nonumber\\
 A_3&=E-XY, & A_4&=c-Yk-\eta,\nonumber\\
 A_5&=\tau^2+(2\tau-k)T-1,&&\nonumber\\
 A_6&=2acS+S^2-Hc^2-1,&&\nonumber\\
 A_7&=k-1-hE-R,&&\nonumber\\
 A_8&=(K_0+4w)C-(q-1)z-F,&&\nonumber\\
 A_9&=V+c-fo, & A_{10}&=V+k-1-cj-hE,\nonumber\\
 A_{11}&=I-(2r+1)c^2, & A_{12}&=I^2-\delta(f^2-1),\nonumber\\
 A_{13}&=I^2(V^2-y^2)+y^2-1, & A_{14}&=B_*.
 \label{eq:int26rows}
\end{align}
The requested single equation is
\begin{equation}
 \boxed{\displaystyle Q_{26}=\sum_{i=1}^{14}A_i^2=0.}
 \label{eq:int26}
\end{equation}
All quantifiers range over $\Z$.  The abbreviations introduce no additional
quantified coordinates and contain only fixed nonnegative powers.

\begin{proposition}
The equation $Q_{26}=0$ has an integer solution exactly when the literal
source-71 quartic has an integer solution.  Its exact total degree is eight.
\end{proposition}
\begin{proof}
The source quartic has two rows $G=\operatorname{Pos}$ and
$Z=\operatorname{Pos}$, each with its own private triple.  Replace these
six coordinates and two rows by $t,p_1,p_2,p_3$ and $B_*=0$.
Lemma~\ref{lem:int-pair} proves equivalence on the projection to all the
other coordinates.  This gives a 36-variable system of 24 residuals.

Translate the first root by
\[
 \tau_{\rm old}=\tau+E(c-\eta).
\]
Then eliminate the following ten monically defined coordinates:
\[
\begin{array}{c|c@{\qquad}c|c}
 X&2u^2 & C_2&c^2\\
 F_2&f^2 & B_{\rm aux}&I^2\\
 g&y^2-V^2 & \Delta&a^2+4a+3\\
 D&2u^2+ac+\gamma(4a+3)&P&k+E(c-\eta)\\
 L&q^2-1&N&q^3
\end{array}
\]
These are integer-polynomial restorations with integer-polynomial inverses
on the zero sets, given by forgetting the defined coordinates and undoing
the root translation.  The ten definition rows become identities.
For the first and main norms the exact identities are
\begin{align*}
 (\tau+T)^2-(k+T)^2+(k+T)k-1
   &=\tau^2+(2\tau-k)T-1,\\
 (ac+S)^2-(a^2+H)c^2-1
   &=2acS+S^2-Hc^2-1.
\end{align*}
The surviving rows are exactly \eqref{eq:int26rows}, up to multiplication
of individual residuals by $-1$, which leaves their squares unchanged.
The literal source-row restorations are also checked by exact sparse
polynomial arithmetic in the archive.

Every residual has degree at most four.  The only residual containing
$t$ is $B_*$, whose $t^4$ coefficient is $-1$.  Therefore
$[t^8]Q_{26}=1$, independently of the fixed coefficients.  This proves
exact degree eight on the actual coefficient slice.
\end{proof}

\subsection{The degree-ten equation in 24 integers}

Retain the same definitions, but now set
\begin{equation}
 c=Yk+\eta,\qquad V=fo-c,
 \qquad y_{\rm old}=V+\nu.
 \label{eq:int24changes}
\end{equation}
Here $c,V$ are polynomial abbreviations.  Replace the coordinate $y$ by
$\nu$ and remove the coordinates $c,V$.  The 24 independent coordinates are
\[
\begin{gathered}
 J,F,\alpha,z,f,h,r,v,w,\tau,\eta,\zeta,\nu,\gamma,o,j,\\
 E,a,I,u,t,p_1,p_2,p_3.
\end{gathered}
\]
Remove $A_4,A_9$, which now vanish identically.  Replace $A_{13}$ by
\begin{equation}
 A'_{13}=(V+\nu)^2-I^2\nu(2V+\nu)-1.
 \label{eq:int24aux}
\end{equation}
In every other residual substitute the two definitions in
\eqref{eq:int24changes}.  The complete single equation is
\begin{equation}
 \boxed{\begin{aligned}
 Q_{24}={}&A_1^2+A_2^2+A_3^2+A_5^2+A_6^2+A_7^2\\
         &+A_8^2+A_{10}^2+A_{11}^2+A_{12}^2+(A'_{13})^2+B_*^2=0.
 \end{aligned}}
 \label{eq:int24}
\end{equation}

\begin{theorem}[Two root translations and two coordinate eliminations]
The integer zero sets of $Q_{26}$ and $Q_{24}$ are related by inverse
integer-polynomial maps.  The polynomial $Q_{24}$ has exactly 24 live
integer variables and exact total degree ten.
\end{theorem}
\begin{proof}
The two definition rows $A_4,A_9$ are monic in $c,V$.  Their right-hand
sides in \eqref{eq:int24changes} use only retained coordinates and have
degree two.  The auxiliary-root translation has the integral inverse
$\nu=y_{\rm old}-V$.  The exact identity
\[
 I^2\bigl(V^2-(V+\nu)^2\bigr)+(V+\nu)^2-1
  =(V+\nu)^2-I^2\nu(2V+\nu)-1
\]
proves restoration of the auxiliary norm for arbitrary signed integers.
There are no sign assumptions or divisibility conditions in these maps.

The degree savings require both root translations.  After eliminating
$c$, the expression $T=E(c-\eta)$ has degree three.  Centering the first
root makes its norm $\tau^2+(2\tau-k)T-1$ of degree four instead of six.
After eliminating $V$, the auxiliary-root translation in
\eqref{eq:int24aux} makes its norm degree five instead of six.  In the
main norm the cancellation of $a^2c^2$ already displayed leaves
$2acS+S^2-Hc^2-1$, also of degree five.  Finally
$I-(2r+1)c^2$ has degree five.  All other residuals have degree at most
four.  In the order displayed in \eqref{eq:int24}, their degrees are
\[
 2,4,3,4,5,4,2,3,5,4,5,4.
\]
Thus $\deg Q_{24}\le10$.

Only $A_{11}$ contains $r$.  The top part of
$c=(a-E)(\eta+\zeta)+\eta$ gives
\[
 [r^2a^4\eta^4]Q_{24}=4.
\]
This coefficient is independent of all six fixed numerals, so the degree
is exactly ten after actual coefficient specialization.  Exact symbolic
checks also verify that every listed coordinate occurs.
\end{proof}

\subsection{An optional 23-variable equation of degree sixteen}

Starting from $Q_{24}$, substitute simultaneously
\begin{equation}
 u_{\rm old}=2q\,u,\qquad w=2q\,u^2.
 \label{eq:int23dyadic}
\end{equation}
Then $X=8q^2u^2$ and $A_1=X-4qw$ is identically zero.  Remove $w$ from
the coordinate list and remove this zero residual.  The other eleven
squared residuals define $Q_{23,16}$, a polynomial in 23 unrestricted
integer variables, with the same six fixed coefficients.

\begin{proposition}
The absence of integer zeros of $Q_{23,16}$ has the same consistency
equivalence as $Q_{24}$.  Its exact degree is sixteen.
\end{proposition}
\begin{proof}
Any new zero restores a $Q_{24}$ zero by the integral substitutions
\eqref{eq:int23dyadic}.  Completeness uses fresh canonical native witnesses,
rather than asserting that every old zero satisfies the restriction.
Source 64's canonical witnesses have
\[
 q=2^n,\quad n\ge4,\qquad R\text{ odd},\quad R>3q,
 \qquad u_{\rm old}=2^{(R-1)/2}.
\]
Since $3\cdot2^n>2n+3$ for $n\ge4$, the exponent
$(R-1)/2-n-1$ is a nonnegative integer.  Thus $2q$ divides
$u_{\rm old}$, and the new $u$ is integral.  The old equation
$2u_{\rm old}^2=4qw$ then gives exactly $w=2qu^2$.
The source-71 determinant construction retains those canonical coordinates,
and the new joint guard and root translations do not change this fact.

With $X=8q^2u^2$, the main-norm residual has degree eight.  Every other
residual has degree at most five.  Its unique top contribution
$X^2=64q^4u^4$ gives
\[
 [J^8u^8]Q_{23,16}=4096\beta^8>0.
\]
Therefore the exact degree is sixteen.  This step is an existence
equivalence obtained through canonical witnesses, not a bijection of all
old zero sets.
\end{proof}

\subsection{Consistency theorem and verification boundary}

Composing the proved transformations with the named source-71 correctness
theorem yields, by ordinary arguments formalizable in ZFC, for each of
$Q_{26},Q_{24},Q_{23,16}$,
\[
 \mathrm{ZFC}\vdash
 \bigl(\neg\exists\mathbf{x}\in\Z^m\;Q(\mathbf{x})=0\bigr)
 \;\longleftrightarrow\;\operatorname{Con}(\mathrm{ZFC}).
\]
The consistency predicate is the specified syntactic ZFC refutation
predicate of the inherited calibration.  The implication does not assert
either consistency or inconsistency of ZFC.  No variable coefficient,
variable exponent, hidden inequality, or quantified abbreviation is used.

The release checker reconstructs all six actual fixed integers, checks
their complete raw hashes, verifies every restored residual exactly,
expands each sum of squares over the coefficient ports, proves the
reported degree by the displayed nonzero coefficient, and checks variable
liveness.  Its finite signed and nearest-square samples are only regression
diagnostics; Lemma~\ref{lem:int-pair} is the universal proof.  The inherited
unbounded native compiler, Pell classification and rank arguments, and
literal machine/refutation calibration are not replaced by those finite
checks.  This package adds no end-to-end proof-assistant or internally
encoded ZFC certificate.
